% One-page letter / statement template. Matches the CV letterhead.
\documentclass[11pt, letterpaper]{article}
\usepackage[T1]{fontenc}
\usepackage[utf8]{inputenc}
\usepackage{charter}
\usepackage[top=0.55in, bottom=0.45in, left=0.8in, right=0.8in]{geometry}
\usepackage[dvipsnames]{xcolor}
\definecolor{linkblue}{HTML}{1155CC}
\usepackage[colorlinks=true, urlcolor=linkblue, linkcolor=linkblue]{hyperref}
\usepackage{parskip}
\pagestyle{empty}
\setlength{\parindent}{0pt}
\linespread{1.0}

\newcommand{\letterhead}{%
  \begin{center}
    {\large\bfseries Rodrigo França Chaves}\\[2pt]
    {\small
      Chicago, IL \enspace$\cdot$\enspace
      \href{mailto:rchaves@uchicago.edu}{rchaves@uchicago.edu} \enspace$\cdot$\enspace
      +1 312 217 5526 \enspace$\cdot$\enspace
      \href{https://rodrigofch7.github.io}{rodrigofch7.github.io}
    }
  \end{center}
  \vspace{2pt}
  {\rule{\linewidth}{0.8pt}}
  \vspace{10pt}
}

\begin{document}
\letterhead
\begin{center}{\bfseries\large Statement of Interest}\\[2pt]
{\small EPIC Faculty Fellowship, Academic Year 2026--2027}\end{center}
\vspace{6pt}

I work on how regulation shapes the physical systems people depend on, and on whether the evidence we use to judge that regulation actually supports the claims made from it. My route into energy came through utilities. My first-authored article in \textit{Utilities Policy} asks whether Brazil's transfer of water and wastewater systems to private operators improved public health, comparing propensity-matched municipalities before and after transfer between 1998 and 2021. It finds significant reductions in diarrheal morbidity among children under five, no effect on diseases unrelated to sanitation, and considerable variation across municipalities. Writing it taught me that the interesting part of a regulated-industry question is usually not the headline average but where it fails to hold.

Two strands of my work since then point directly at EPIC's agenda. The first is the local cost of energy infrastructure. In a recent project I assembled a dataset of Chicago-area data centers, confirming and geocoding 45 facilities from 124 listed sites, dated each opening from air permit filings between 2000 and 2024, and matched those dates to ZIP-level home values and to American Community Survey household electricity and water cost data across 660 metropolitan ZIP codes. The design is descriptive rather than causal, and I was explicit about that, but it convinced me that the distributional question of who pays for energy buildout is badly under-measured and increasingly urgent.

The second strand is measurement itself. At the World Bank this summer I estimated how construction responded to a new bus rapid transit corridor in Dakar using difference-in-differences over 100-meter grid cells built from Google Open Buildings and GHSL satellite rasters. The most useful thing I did there was not the estimate. It was noticing that the project's building-detection threshold had been set at nearly twice the value its source paper recommends for that region, which discarded real structures most heavily in dense and informal neighborhoods. Because that bias varies by neighborhood type, it would not have differenced out of the design. Environmental and energy questions lean heavily on remote sensing and on administrative records collected for other purposes, and I have learned to treat the measurement layer as part of the identification problem rather than as a preliminary.

My goal for this fellowship is to move from adjacent work into energy and environmental economics proper, under supervision, on questions where the answer has a decision attached to it. I am drawn in particular to the enforcement-targeting agenda: regulators cannot inspect everything, and the choice of what to inspect is an empirical question with real distributional consequences. That sits at the intersection of what I already do, which is prediction over large administrative data, and what I want to learn, which is how those tools hold up when they carry regulatory weight. A second thread interests me as much: constructing datasets that do not yet exist. I have built pipelines that turn century-old cloth-bound registers into structured records linked to historical censuses, and gridded panels assembled from satellite rasters, and in both cases the measurement layer determined which questions could be asked at all. I would also value the opportunity to produce something citable. My strongest work so far has come from projects where I was responsible for a result end to end, and I would like this fellowship to end in a paper or report I can defend.

Beyond the research, I intend to build a career at the boundary between economic evidence and regulatory practice, whether in an institution like the World Bank, in economic consulting, or in policy research. EPIC is the place at Chicago where that boundary is taken seriously as a research problem in its own right.

\end{document}
